\chapter*{Задачи за упражнение}
\addcontentsline{toc}{chapter}{Задачи за упражнение}
\markboth{Задачи за упражнение}{Задачи за упражнение}

\begin{rembox}[Как да ползвате тази глава]
Задачите по-долу са подредени по темите на учебника - за всяка тема
има по три теоретични задачи, решими само с~материала до и~включително
съответната тема (без да се налага да прелиствате напред). Решенията
са събрани в~отделна глава в~самия край на книгата, за да можете
първо да опитате сами. Означенията навсякъде съвпадат с~тези в~съответната тема.
\end{rembox}

\section*{Тема 1: Въведение}
\addcontentsline{toc}{section}{Тема 1: Въведение}

\begin{probbox}[Задача 1.1]
Нека $\Omega$ е разбито на $k$ на брой непразни, по двойки
несъвместими блокове $B_1,\dots,B_k$ с~$B_1\cup\dots\cup B_k=\Omega$.
Нека $\mathcal A=\sigma(B_1,\dots,B_k)$ е породената от тях
$\sigma$-алгебра. Докажете, че $|\mathcal A|=2^k$.
\end{probbox}

\begin{probbox}[Задача 1.2]
Нека $\{\mathcal A_i\}_{i\in I}$ е произволно (дори неизброимо)
семейство от $\sigma$-алгебри върху едно и~също $\Omega$. Докажете,
че $\mathcal A:=\bigcap_{i\in I}\mathcal A_i$ също е $\sigma$-алгебра
върху $\Omega$.
\end{probbox}

\begin{probbox}[Задача 1.3]
Нека $\mathcal A$ е $\sigma$-алгебра върху $\Omega$ и~$A,B\in
\mathcal A$. Докажете директно от аксиомите (без да ползвате готови
свойства от Следствие 1.1), че $A\setminus B\in\mathcal A$ и~$A\triangle B:=(A\setminus B)\cup(B\setminus A)\in\mathcal A$.
\end{probbox}

\section*{Тема 2: Вероятност и~вероятностно пространство}
\addcontentsline{toc}{section}{Тема 2: Вероятност и~вероятностно пространство}

\begin{probbox}[Задача 2.1]
Докажете неравенството на Бонферони: за произволни събития $A,B$,
\[
  \mathbb P(A\cap B) \ge \mathbb P(A)+\mathbb P(B)-1.
\]
\end{probbox}

\begin{probbox}[Задача 2.2]
Докажете чрез индукция по $n$ крайната субадитивност на вероятността:
за произволни събития $A_1,\dots,A_n$,
\[
  \mathbb P\Big(\bigcup_{i=1}^n A_i\Big) \le \sum_{i=1}^n \mathbb P(A_i).
\]
(Използвайте само аксиомите и~вече доказаното свойство за две
събития.)
\end{probbox}

\begin{probbox}[Задача 2.3]
Нека $(A_n)_{n\ge1}$ е редица от събития с~$\mathbb P(A_n)=1$ за
всяко $n$. Докажете, че $\mathbb P\big(\bigcap_{n=1}^{\infty}
A_n\big)=1$.
\end{probbox}

\section*{Тема 3: Равномерна, условна и~пълна вероятност}
\addcontentsline{toc}{section}{Тема 3: Равномерна, условна и~пълна вероятност}

\begin{probbox}[Задача 3.1]
Хвърляме два правилни зара. Намерете вероятността сборът на числата
да е~7, при условие че първият зар е паднал на четно число.
\end{probbox}

\begin{probbox}[Задача 3.2]
Нека $\mathbb P(A)>0$ е фиксирано. Докажете, че функцията
$Q(\cdot):=\mathbb P(\cdot\mid A)$, дефинирана върху $\mathcal A$,
сама по себе си е вероятностна мярка върху $(\Omega,\mathcal A)$,
т.е.\ проверете трите аксиоми: $Q(B)\ge0$ за всяко $B\in\mathcal A$;
$Q(\Omega)=1$; и~$\sigma$-адитивност.
\end{probbox}

\begin{probbox}[Задача 3.3]
Три машини произвеждат съответно $30\%$, $45\%$ и~$25\%$ от всички
детайли в~един цех, с~процент на брак $2\%$, $3\%$ и~$5\%$
съответно. Избираме случаен детайл и~се оказва бракуван. Каква е
вероятността той да е произведен от втората машина?
\end{probbox}

\section*{Тема 4: Случайни величини}
\addcontentsline{toc}{section}{Тема 4: Случайни величини}

\begin{probbox}[Задача 4.1]
Нека $X,Y$ са случайни величини във $V=(\Omega,\mathcal A,\mathbb
P)$. Докажете, че $Z=\max(X,Y)$ също е случайна величина.
\end{probbox}

\begin{probbox}[Задача 4.2]
Нека $A,B\in\mathcal A$. Използвайки свойствата на индикаторната
функция, докажете, че
\[
  \mathbb P(A\cup B) = \mathbb E I_A + \mathbb E I_B - \mathbb E(I_A I_B).
\]
\end{probbox}

\begin{probbox}[Задача 4.3]
Нека $X,Y$ са независими дискретни случайни величини и~$g,h:\mathbb
R\to\mathbb R$ са произволни функции. Докажете, че $g(X)$ и~$h(Y)$
също са независими случайни величини.
\end{probbox}

\section*{Тема 5: Очакване, дисперсия, пораждащи функции}
\addcontentsline{toc}{section}{Тема 5: Очакване, дисперсия, пораждащи функции}

\begin{probbox}[Задача 5.1]
Докажете, че $DX\ge0$ за всяка случайна величина $X$ с~$\mathbb
EX^2<\infty$, и~че $D(aX+b)=a^2DX$ за произволни константи
$a,b\in\mathbb R$.
\end{probbox}

\begin{probbox}[Задача 5.2]
Нека $X\sim Bin(n,p)$ с~пораждаща функция $g_X(s)=(ps+1-p)^n$.
Използвайки формулите $g_X'(1)=\mathbb EX$ и~$DX=g_X''(1)+g_X'(1)-
(g_X'(1))^2$, пресметнете $\mathbb EX$ и~$DX$ директно от $g_X$, и~сравнете с~познатите стойности $np$, $np(1-p)$.
\end{probbox}

\begin{probbox}[Задача 5.3]
Намерете пораждащата функция $g_X(s)$ на $X\sim Bin(n,p)$, тръгвайки
директно от дефиницията $g_X(s)=\sum_k s^k\,\mathbb P(X=k)$ и~биномната формула на Нютон (т.е.\ изведете формулата от Задача~5.2,
вместо да я приемате наготово).
\end{probbox}

\section*{Тема 6: Основни дискретни разпределения}
\addcontentsline{toc}{section}{Тема 6: Основни дискретни разпределения}

\begin{probbox}[Задача 6.1]
(Липса на памет на геометричното разпределение) Нека $X\sim Geom(p)$
с~носител $\{1,2,3,\dots\}$. Докажете, че за всички цели $m,n\ge0$:
\[
  \mathbb P(X>m+n \mid X>m) = \mathbb P(X>n).
\]
\end{probbox}

\begin{probbox}[Задача 6.2]
Нека $X\sim Bin(n,p)$ и~$Y\sim Bin(m,p)$ са независими. Докажете, че
$X+Y\sim Bin(n+m,p)$.
\end{probbox}

\begin{probbox}[Задача 6.3]
Нека $X\sim Bin(1000,\,0.002)$. Използвайки теоремата на Поасон за
сходимост по разпределение, оценете приблизително $\mathbb P(X=3)$
чрез подходящо разпределение на Поасон, и~сравнете качествено
(без калкулатор) с~точната биномна вероятност.
\end{probbox}

\section*{Тема 7: Ковариация, корелация}
\addcontentsline{toc}{section}{Тема 7: Ковариация, корелация}

\begin{probbox}[Задача 7.1]
Докажете, че за произволни случайни величини $X,Y$ с~крайни
дисперсии и~константи $a,b,c,d\in\mathbb R$:
\[
  \operatorname{cov}(aX+b,\,cY+d) = ac\cdot\operatorname{cov}(X,Y).
\]
\end{probbox}

\begin{probbox}[Задача 7.2]
Докажете, че $D(X+Y) = DX+DY+2\operatorname{cov}(X,Y)$ за произволни
случайни величини $X,Y$ с~крайни дисперсии.
\end{probbox}

\begin{probbox}[Задача 7.3]
Нека $X$ е равномерно разпределена върху $\{-1,0,1\}$ (всяка стойност
с~вероятност $1/3$) и~$Y=X^2$. Докажете, че $\operatorname{cov}(X,Y)=0$,
но $X$ и~$Y$ \textbf{не} са независими.
\end{probbox}

\section*{Тема 8: Условно математическо очакване}
\addcontentsline{toc}{section}{Тема 8: Условно математическо очакване}

\begin{probbox}[Задача 8.1]
Използвайки кула-свойството и~линейността на условното очакване,
докажете, че
\[
  \mathbb E\big[X\cdot \mathbb E[X\mid Y]\big] = \mathbb E\big[(\mathbb E[X\mid Y])^2\big].
\]
\end{probbox}

\begin{probbox}[Задача 8.2]
Хвърляме правилен зар, за да получим $N\in\{1,\dots,6\}$, а~след
това хвърляме правилна монета $N$ пъти. Нека $X$ е броят на паднали
се ези. Намерете $\mathbb EX$, като първо пресметнете $\mathbb E[X
\mid N]$.
\end{probbox}

\begin{probbox}[Задача 8.3]
Нека съвместното разпределение на дискретните $X,Y$ е зададено с~$\mathbb P(X=0,Y=0)=0.1$, $\mathbb P(X=0,Y=1)=0.2$, $\mathbb
P(X=1,Y=0)=0.3$, $\mathbb P(X=1,Y=1)=0.4$. Намерете $\mathbb E[X\mid
Y=0]$ и~$\mathbb E[X\mid Y=1]$.
\end{probbox}

\section*{Тема 9: Непрекъснати случайни величини}
\addcontentsline{toc}{section}{Тема 9: Непрекъснати случайни величини}

\begin{probbox}[Задача 9.1]
Нека $f_X(x)=c(1-x^2)$ за $x\in(-1,1)$ и~$f_X(x)=0$ извън този
интервал. Намерете $c$, функцията на разпределение $F_X$, и~$\mathbb P(0<X<1/2)$.
\end{probbox}

\begin{probbox}[Задача 9.2]
Нека $X\sim U(0,1)$ и~$Y=-\ln X$. Използвайки теоремата за смяна на
променливите (Тема~9), докажете, че $Y\sim Exp(1)$.
\end{probbox}

\begin{probbox}[Задача 9.3]
Нека $X\ge0$ е непрекъсната случайна величина с~плътност $f_X$.
Докажете, че
\[
  \mathbb EX = \int_0^{\infty} \mathbb P(X>x)\,dx.
\]
\emph{(Упътване: разпишете дясната страна като двоен интеграл и~разменете реда на интегриране.)}
\end{probbox}

\section*{Тема 10: Видове непрекъснати случайни величини}
\addcontentsline{toc}{section}{Тема 10: Видове непрекъснати случайни величини}

\begin{probbox}[Задача 10.1]
(Липса на памет на експоненциалното разпределение) Нека $X\sim
Exp(\lambda)$. Докажете, че за всички $s,t\ge0$:
\[
  \mathbb P(X>s+t\mid X>s) = \mathbb P(X>t).
\]
\end{probbox}

\begin{probbox}[Задача 10.2]
Нека $X\sim N(\mu,\sigma^2)$ и~$Z=\dfrac{X-\mu}\sigma$. Използвайки
теоремата за смяна на променливите (Тема~9), докажете, че $Z\sim
N(0,1)$.
\end{probbox}

\begin{probbox}[Задача 10.3]
Намерете медианата на $X\sim Exp(\lambda)$, т.е.\ стойността $m$,
за която $\mathbb P(X\le m)=1/2$, изразена чрез $\lambda$.
\end{probbox}

\section*{Тема 11: Непрекъснати съвместни разпределения, Гама разпределение}
\addcontentsline{toc}{section}{Тема 11: Непрекъснати съвместни разпределения, Гама разпределение}

\begin{probbox}[Задача 11.1]
Нека $(X,Y)$ е равномерно разпределена върху единичния квадрат
$(0,1)\times(0,1)$. Намерете $\mathbb P(X+Y<1)$.
\end{probbox}

\begin{probbox}[Задача 11.2]
Нека $X_1,\dots,X_n$ са независими, всяка $X_i\sim Exp(\lambda)$.
Използвайки, че $Exp(\lambda)=\Gamma(1,\lambda)$ и~адитивността на
Гама разпределението (Тема~11), докажете, че $\sum_{i=1}^n X_i \sim
\Gamma(n,\lambda)$.
\end{probbox}

\begin{probbox}[Задача 11.3]
Нека $(X,Y)$ има съвместна плътност $f(x,y)=x+y$ за $(x,y)\in(0,1)^2$
(и~$0$ извън квадрата). Намерете $\operatorname{cov}(X,Y)$.
\end{probbox}

\section*{Тема 12: Сходимост и~закони за големите числа}
\addcontentsline{toc}{section}{Тема 12: Сходимост и~закони за големите числа}

\begin{probbox}[Задача 12.1]
Нека $X_n\xrightarrow{\mathbb P}c$ за константа $c\in\mathbb R$.
Докажете директно от дефиницията на сходимост по вероятност, че
$X_n^2\xrightarrow{\mathbb P}c^2$.
\end{probbox}

\begin{probbox}[Задача 12.2]
Нека $X_1,\dots,X_n$ са i.i.d.\ с~$DX_1=\sigma^2=4$. Използвайки
неравенството на Чебишов (както в~доказателството на СЗГЧ), намерете
най-малкото $n$, за което е гарантирано $\mathbb P\big(|\overline
X_n-\mu|>0.5\big)\le0.05$.
\end{probbox}

\begin{probbox}[Задача 12.3]
Хвърляме монета с~неизвестна вероятност $p$ за ези $n=10000$ пъти и~означаваме с~$\hat p_n$ дела на паднали се ези. Използвайки
неравенството на Чебишов и~оценката $DX_i\le\tfrac14$ за произволна
Бернулиева величина, оценете горна граница за $\mathbb P(|\hat p_n-
p|>0.02)$.
\end{probbox}

\section*{Тема 13: Централна гранична теорема}
\addcontentsline{toc}{section}{Тема 13: Централна гранична теорема}

\begin{probbox}[Задача 13.1]
Хвърляме правилна монета $n=100$ пъти и~означаваме с~$S_n$ броя на
падналите се ези. Използвайки ЦГТ (респ.\ следствието на
Моавър--Лаплас), оценете приблизително $\mathbb P(45<S_n<55)$ чрез
$\Phi$.
\end{probbox}

\begin{probbox}[Задача 13.2]
Нека $X\sim Exp(\lambda)$ с~пораждаща функция на моментите
$M_X(t)=\dfrac\lambda{\lambda-t}$ за $t<\lambda$. Използвайки
свойство б) от Тема~13 ($M_X^{(k)}(0)=\mathbb EX^k$), пресметнете
$\mathbb EX$ и~$DX$ директно от $M_X$.
\end{probbox}

\begin{probbox}[Задача 13.3]
Нека $X_i\sim Be(0.5)$ i.i.d. Пресметнете границата, която дава
неравенството на Бери--Есен за $N=100$.
\end{probbox}

\section*{Тема 14: Точкови оценки}
\addcontentsline{toc}{section}{Тема 14: Точкови оценки}

\begin{probbox}[Задача 14.1]
Нека $X\sim Poisson(\lambda)$ и~$\bar X=(X_1,\dots,X_n)$ е извадка от
независими наблюдения. Намерете МПО оценката $\hat\lambda$.
\end{probbox}

\begin{probbox}[Задача 14.2]
Нека $X\sim Exp(\lambda)$. Намерете оценката по метода на моментите
$\hat\lambda$.
\end{probbox}

\begin{probbox}[Задача 14.3]
Нека $\hat\theta_n$ е неизместена оценка на $\theta$ (за всяко $n$) с~$D\hat\theta_n\to0$ при $n\to\infty$. Докажете, че $\hat\theta_n$ е
състоятелна оценка на $\theta$.
\end{probbox}

\section*{Тема 15: Доверителни интервали и~хипотези}
\addcontentsline{toc}{section}{Тема 15: Доверителни интервали и~хипотези}

\begin{probbox}[Задача 15.1]
Нека $X_1,\dots,X_n$ са i.i.d.\ $Be(p)$ и~$\hat p_n=\overline X_n$.
Използвайки ЦГТ, изведете приблизителен (асимптотичен) доверителен
интервал за $p$ с~ниво на доверие $\gamma$.
\end{probbox}

\begin{probbox}[Задача 15.2]
Нека $X\sim N(\mu,4)$ ($\sigma^2=4$ известно) и~извадка от $n=16$
наблюдения дава $\hat\mu=50$. Постройте $95\%$ доверителен интервал
за $\mu$ (използвайте $q_{0.975}\approx1.96$).
\end{probbox}

\begin{probbox}[Задача 15.3]
Нека $X\sim N(\mu,16)$ ($\sigma^2=16$ известно). Проверяваме $H_0:
\mu=50$ срещу $H_1:\mu\neq50$ при ниво на значимост $\alpha=0.05$
($q_{0.975}\approx1.96$). Извадка от $n=25$ наблюдения дава
$\hat\mu=52$. Отхвърля ли се $H_0$?
\end{probbox}

\section*{Тема 16: Линейна регресия}
\addcontentsline{toc}{section}{Тема 16: Линейна регресия}

\begin{probbox}[Задача 16.1]
Дадени са наблюденията $(x_k,y_k)$: $(1,2),(2,3),(3,5),(4,6)$.
Пресметнете оценките $\hat\beta_0,\hat\beta_1$ по метода на
най-малките квадрати.
\end{probbox}

\begin{probbox}[Задача 16.2]
Докажете директно от нормалните уравнения (условията за минимум от
Тема~16), че сборът на остатъците е нула:
\[
  \sum_{k=1}^n\big(Y_k-\hat\beta_0-\hat\beta_1X_k\big) = 0.
\]
\end{probbox}

\begin{probbox}[Задача 16.3]
Докажете „дуалността“ между доверителен интервал и~тест за наклона:
отхвърлянето на $H_0:\beta_1=b$ при ниво на значимост $\alpha$ (по
конструкцията от Тема~16) е еквивалентно на това, че $b$ не попада в~$(1-\alpha)$-доверителния интервал за $\beta_1$.
\end{probbox}
